\section{A square-zero extension and projective dimension}
\label{sec:square-zero}

We first establish the homological statement for an arbitrary
split square-zero extension over a field.  Its derived tensor
powers cannot in general be replaced by ordinary tensor powers.
The decomposition below will let us prove two separate facts: eventual
extinction gives a finite projective resolution, while survival of the
\(r\)th iterate rules out any projective resolution of length less than \(r\).

\subsection{The derived bar decomposition}

The splitting below also follows from the graded-component decomposition
of Minamoto and Yamaura
\cite[Lemma~4.13(4) and the proof of Theorem~4.17]{MinamotoYamaura},
after forgetting the grading and passing to opposite algebras.
The direct proof identifies each derived tensor factor and its shift.

\begin{proposition}
\label{prop:bar-splitting}
Let \(D\) be a finite-dimensional \(k\)-algebra, let \(X\) be
a finite-dimensional \(D\)-bimodule, and let
\(A=D\ltimes X\) be the algebra with multiplication
\[
(d,x)(d',x')=(dd',dx'+xd').
\]
Let \(N\) be a left \(D\)-module, inflated to \(A\) through
the projection \(A\to D\), and put \(F=X\otL_D-\).
In the derived category of left \(D\)-modules there is an
isomorphism
\begin{equation}
\label{eq:bar-splitting}
D\otL_A N
\simeq
\bigoplus_{r\ge0}(F^rN)[r].
\end{equation}
Here \(F^0N=N\), and every \(F^r\) denotes an iteration
of the derived tensor functor.
\end{proposition}

\begin{proof}
Use the ordinary unnormalized bar resolution of \(N\) over
the field \(k\).  It is a free left \(A\)-resolution with
degree \(-n\) term \(A\otimes_k A^{\otimes_k n}\otimes_k N\).
Its augmentation is a quasi-isomorphism: the augmented
bar complex has the usual \(k\)-linear contraction obtained
by inserting the unit.  After tensoring with the right
\(A\)-module \(D\), it yields a model \(\mathcal B\) for
\(D\otL_AN\), with
\[
\mathcal B^{-n}=D\otimes_k A^{\otimes_k n}\otimes_k N
\qquad(n\ge0).
\]
The differential is the alternating sum of the faces
that multiply adjacent algebra entries or apply the two
endpoint actions.

Expand each internal copy of \(A\) as \(D\oplus X\).
Every nonzero face preserves the number of entries from
\(X\).  Indeed, a product involving exactly one \(X\)-entry
remains in \(X\), a product of two \(X\)-entries is zero,
and an endpoint action by \(X\) is zero on both \(D\) and
\(N\).  Thus there is an actual decomposition of complexes
\[
\mathcal B=\bigoplus_{r\ge0}\mathcal B_r,
\]
where \(\mathcal B_r\) consists of words with precisely
\(r\) entries from \(X\).

For fixed \(r\), write \(i_0,\ldots,i_r\) for the lengths
of the blocks of \(D\)-entries before, between, and after
the \(X\)-entries.  The corresponding tensor is
\begin{equation}
\label{eq:bar-block}
D\otimes_kD^{\otimes_k i_0}\otimes_k X
\otimes_kD^{\otimes_k i_1}\otimes_k\cdots
\otimes_k X\otimes_kD^{\otimes_k i_r}\otimes_k N,
\end{equation}
in degree \(-r-\sum_{j=0}^r i_j\).
We identify this with the multiple two-sided bar complex
for
\begin{equation}
\label{eq:derived-multibar}
D\otL_D
\underbrace{X\otL_D\cdots\otL_D X}_{r\ \text{factors}}
\otL_D N,
\end{equation}
shifted by \(r\).

First consider the signs.  Within block \(j\), a bar face
has alternating-sign offset
\[
j+\sum_{h<j}i_h.
\]
In the ordinary multibar total differential, ordered from
left to right, the corresponding offset is
\(\sum_{h<j}i_h\); shifting the total complex by \(r\)
adds the factor \((-1)^r\).
Rescale the component \eqref{eq:bar-block} by
\[
(-1)^{\sum_{j=0}^r(r-j)i_j}.
\]
When a face reduces \(i_j\) by one, this factor changes by
\((-1)^{r-j}\), precisely the discrepancy between the
two differentials.  This gives the required chain
isomorphism.  When \(r\ge1\) and \(i_j=0\), there is no
nonzero face in that block: the prospective face multiplies
two \(X\)-entries or applies an endpoint \(X\)-action.
For \(r=0\), this is the ordinary two-sided bar complex
for \(D\otL_DN\).  At \(i_0=0\) its differential is absent;
the map \(D\otimes_kN\to N\) is the separate augmentation.

We next justify that the multibar computes the genuinely
derived tensor in \eqref{eq:derived-multibar}, without any
flatness assumption on \(X\) over \(D\).
For a bounded-above complex \(Z\) of left \(D\)-modules,
the total left-module bar resolution has terms obtained
from
\[
D\otimes_k D^{\otimes_k i}\otimes_k Z^q
\]
in degree \(q-i\).  Its terms are free left \(D\)-modules
because \(k\) is a field, and its augmentation to \(Z\)
is a quasi-isomorphism.  For each total degree there are
only finitely many possible \(i\), since \(Z\) is bounded
above.  Thus it is a bounded-above complex of free
modules and can be used to compute a derived tensor.

Starting at the right, resolve \(N\) in this way and
tensor with \(X\).  Resolve the resulting left
\(D\)-complex in the same way and tensor with the next
\(X\), continuing for the fixed finite number \(r\)
of factors.  A final bar resolution with the leading
factor \(D\) gives exactly the multibar
\eqref{eq:bar-block}, before its shift.
All exterior \(D\)-actions are preserved.
This establishes its identification with
\eqref{eq:derived-multibar}, which is isomorphic to
\(F^rN\).

Finally, in degree \(-n\) of \(\mathcal B\) only
\(0\le r\le n\) occur, and for each such \(r\) there
are finitely many tuples satisfying
\(r+\sum_j i_j=n\).
All totalizations above therefore use direct sums that
are finite in each degree.  There is no passage to
products or completed tensor powers.
Direct sums are exact for modules, so the componentwise
quasi-isomorphisms give \eqref{eq:bar-splitting}.
The argument uses no connectedness assumption on \(D\);
in particular it applies to the product algebra of
\Cref{sec:ordinary-simulation}.
\end{proof}

\subsection{Detecting finite projective dimension and lower bounds}

\begin{corollary}
\label{cor:pd-detection}
In the setting of \Cref{prop:bar-splitting}, suppose that
\(D\) has finite global dimension on both sides and that
\(N\) is finite-dimensional.
\begin{enumerate}
\item If \(F^tN\simeq0\) for some \(t\ge1\), then
      \(\pd_A N<\infty\).
\item If \(F^rN\not\simeq0\), then \(\pd_A N\ge r\),
      with infinite projective dimension allowed.
\end{enumerate}
\end{corollary}

\begin{proof}
Every fixed iterate \(F^rN\) has bounded,
finite-dimensional cohomology.
For example, if \(g\) bounds the right global dimension
of \(D\), then for every finite-dimensional left
\(D\)-module \(V\) the groups
\(\Tor_i^D(X,V)\) are finite-dimensional and vanish for
\(i>g\).  Applying this observation to successive
truncation triangles shows inductively that
\[
H^q(F^rN)=0\quad\text{unless}\quad -rg\le q\le0.
\]
In particular, all iterates lie in \(\D^{\le0}(D)\).

For the first assertion, vanishing at \(t\) implies
vanishing at every \(r\ge t\).
By \eqref{eq:bar-splitting},
\(Z=D\otL_AN\) then has bounded cohomology.
Since \(X\) is a nilpotent ideal, it is contained in
\(\rad A\), and
\[
\rad A=(\rad D)\oplus X,
\qquad A/\rad A\cong D/\rad D=:S.
\]
Derived associativity gives
\[
S\otL_A N
\simeq S\otL_D(D\otL_A N)
=S\otL_D Z.
\]
The right \(D\)-module \(S\) has a finite projective
resolution, so this complex also has bounded cohomology.

Choose a minimal projective resolution of \(N\) over
the finite-dimensional algebra \(A\), using successive
projective covers.  Its terms are finitely generated,
and every differential has image in the radical of
its target.  Tensoring with \(S=A/\rad A\) therefore
annihilates every differential.  If \(Q\) is a nonzero
finitely generated projective \(A\)-module, then
\[
S\otimes_A Q\cong Q/(\rad A)Q\ne0
\]
by Nakayama's Lemma.  Consequently bounded cohomology
of \(S\otL_AN\) forces the minimal resolution to have
only finitely many nonzero terms.  This proves
\(\pd_A N<\infty\).

For the second assertion, choose \(q\le0\) with
\(H^q(F^rN)\ne0\).  Its direct summand in
\eqref{eq:bar-splitting} gives
\[
H^{q-r}(D\otL_A N)\ne0,\qquad q-r\le-r.
\]
A projective \(A\)-resolution of \(N\) of length
strictly less than \(r\) would make the derived
tensor zero in this degree.  Hence
\(\pd_A N\ge r\), as asserted.
\end{proof}

\begin{remark}
For inflated modules, the projective-dimension formula
of Minamoto and Yamaura
\cite[Corollary~4.11]{MinamotoYamaura}, after passage
to the opposite algebra to match our left-module
convention, gives the stronger equality
\[
\pd_A N
=\sup_{r\ge0}\bigl\{\pd_D(F^rN)+r\bigr\},
\]
where projective dimension on the right is understood
for derived complexes, with \(\pd_D(0)=-\infty\).
The graded and ungraded dimensions
agree by \cite[Proposition~2.5 and Lemma~3.6]{MinamotoYamaura}.
Their
\cite[Theorem~4.17]{MinamotoYamaura} also gives the
corresponding perfectness criterion.
The direct bar and minimal-resolution arguments above
prove the assertions needed here.
\end{remark}

\subsection{The fixed algebra and its modules}

\begin{proof}[Proof of \Cref{thm:main}]
Work over \(k=\C\).
Fix the algebra \(R\), its central idempotent \(e\),
the automorphism \(\alpha\), and the selection functor
\(H\) supplied by \Cref{prop:selection}.
Apply \Cref{thm:realization} to these fixed data,
obtaining \(B\) and \(P\) satisfying
\eqref{eq:realization}.
Choose the support interval \([a_*,b_*]\) once, and
construct \(C,W,O,D,X\) as in
\Cref{sec:ordinary-simulation}.
Finally fix the ordinary finite-dimensional unital
algebra
\[
A=D\ltimes X.
\]
These choices are made independently of \(m\), as summarized in
\Cref{tab:fixed-varying}.

\begin{table}[htbp]
\centering
\begin{tabular}{@{}p{0.24\textwidth}p{0.68\textwidth}@{}}
\toprule
Choice & Objects \\
\midrule
Fixed once & \(R,e,\alpha\); the presentation and algebra \(B\);
the lifted diagram and complex \(P\); the interval \([a_*,b_*]\);
\(C,W,O,T,D,X\); and \(A=D\ltimes X\). \\
Varying with \(m\) & The quotient \(F_m\), its central character,
the selection module \(Y_m\), and the inflated module
\(N_m=M(Y_m)\). \\
\bottomrule
\end{tabular}
\caption{The quantifier order in the construction of \(A\). Varying the
test module changes neither the tensor functor nor the final algebra.}
\label{tab:fixed-varying}
\end{table}

For a finite-dimensional left \(R\)-module \(Y\),
regard \(M(Y)\) as supported on the first factor
of \(D\).  By \eqref{eq:compression} and
\eqref{eq:realization},
\[
F^2M(Y)
\simeq M(H(Y))[b_*]\oplus M(H(Y))[b_*+3].
\]
Induction, using preservation of finite direct sums
and shifts by derived tensor, gives
\begin{equation}
\label{eq:even-iterates}
F^{2j}M(Y)
\simeq
\bigoplus_{q=0}^{j}
M(H^{\circ j}(Y))[jb_*+3q]^{\oplus\binom{j}{q}}
\qquad(j\ge0).
\end{equation}
The induction applies the realization isomorphism
separately to each module \(H^{\circ j}(Y)\); it requires
no splitting natural in \(Y\).

For each \(m\ge1\), let \(Y_m\) be the finite-dimensional
module supplied by \eqref{eq:extinction}, and let
\(N_m=M(Y_m)\), first as a left \(D\)-module on its
first factor and then as an inflated left \(A\)-module.
It is finite-dimensional.
Since \(H^{\circ m}(Y_m)=0\) and \(H^{\circ(m-1)}(Y_m)\ne0\),
\eqref{eq:even-iterates} yields
\[
F^{2m}N_m\simeq0,\qquad
F^{2m-2}N_m\not\simeq0.
\]
The latter nonvanishing follows because a finite
direct sum of shifts of the nonzero module
\(M(H^{\circ(m-1)}(Y_m))\) is nonzero. Shifts preserve both extinction
and nonvanishing; the lower bound now uses the iterate index
\(r=2m-2\) in \Cref{cor:pd-detection}, not a new estimate on the shifts.
The algebra \(D\) has finite global dimension on both
sides by \Cref{prop:ordinary-simulation}.
Applying \Cref{cor:pd-detection} proves
\[
2m-2\le\pd_A N_m<\infty
\qquad(m\ge1).
\]
All \(N_m\) are finitely generated over the single
fixed algebra \(A\), while these lower bounds are
unbounded.  Therefore \(\findim A=\infty\), which
refutes the asserted finiteness for all
finite-dimensional algebras.
\end{proof}

\subsection{The auxiliary left--right asymmetry}
\label{sec:left-right-asymmetry}

\begin{proof}[Proof of \Cref{cor:left-right-asymmetry}]
Let \(A\) be the algebra of \Cref{thm:main}, and choose a basic
finite-dimensional complex algebra \(C\) Morita equivalent to
\(A^{\mathrm{op}}\). The right-module
little and big finitistic dimensions of \(C\) are infinite, and its
injective right modules do not generate: these are the conclusions
for left \(A\)-modules in \Cref{sec:introduction}, transported through
Morita equivalence. Indeed, this equivalence preserves projective dimensions
and finite generation, and the induced equivalence of unbounded derived
categories preserves injectives and coproducts.

Choose primitive idempotents \(e_1,\ldots,e_n\) representing the simple
right \(C\)-modules \(S_C(e_i)=e_iC/e_i\rad C\). In Cummings'
Construction~2.1 \cite{Cummings2024}, set
\[
E=\prod_{i=1}^n\C[\varepsilon_i]/(\varepsilon_i^2),
\qquad U=\bigoplus_{i=1}^n S_C(e_i),
\qquad
T=\begin{pmatrix}C&0\\ {}_E U_C&E\end{pmatrix}.
\]
The left \(E\)-action on \(U\) factors through
\(E/\rad E\cong\C^n\), with the \(i\)th factor acting on \(S_C(e_i)\).
This is a finite-dimensional complex algebra.

Write \(\findim_r\) and \(\Findim_r\) for the right-module dimensions,
to display Cummings' convention. By
\cite[Propositions~3.2--3.3]{Cummings2024},
\[
\begin{gathered}
\findim_r C\le\findim_r T,\\
\Findim_r C\le\Findim_r T,\\
\findim_r(T^{\mathrm{op}})=\Findim_r(T^{\mathrm{op}})=0.
\end{gathered}
\]
His Proposition~4.3 says that injective generation for right
\(T\)-modules implies it for right \(C\)-modules, while
Proposition~4.4 establishes it for right \(T^{\mathrm{op}}\)-modules
\cite[Propositions~4.3--4.4]{Cummings2024}. Hence both right finitistic
dimensions of \(T\) are infinite and its injective right modules do not
generate; for
\(T^{\mathrm{op}}\), both right dimensions are zero and its injective
right modules generate. Taking \(\Lambda=T^{\mathrm{op}}\) translates
these statements into our left-module convention.
\end{proof}
